% Generated draft: requires mathematical and source review.
\documentclass[11pt]{article}
\usepackage[margin=25mm]{geometry}
\usepackage{fontspec}
\setmainfont{Latin Modern Roman}
\usepackage{amsmath,amssymb,mathtools,unicode-math}
\setmathfont{Latin Modern Math}
\usepackage{enumitem,needspace,xurl,hyperref}
\setlength{\parindent}{0pt}
\setlength{\parskip}{5pt}
\setlength{\emergencystretch}{3em}
\newcommand{\N}{\mathbb{N}}
\newcommand{\Z}{\mathbb{Z}}
\newcommand{\Q}{\mathbb{Q}}
\newcommand{\R}{\mathbb{R}}
\newcommand{\C}{\mathbb{C}}
\newcommand{\F}{\mathbb{F}}
\newcommand{\A}{\mathbb{A}}
\newcommand{\PP}{\mathbb{P}}
\newcommand{\E}{\mathbb{E}}
\newcommand{\eps}{\varepsilon}
\newcommand{\dd}{\,\mathrm{d}}
\begin{document}
\section*{UnsolvedMath ID: 800044}
\textbf{Source problem code:} catalogue entry
\subsection{Definitions and mathematical statement}
A semigroup is a nonempty set equipped with an associative binary operation; an identity element is not required. A variety of semigroups is the class of all semigroups satisfying some fixed set of identities in the language of one binary operation. For a semigroup variety \(\mathcal V\), define
\[
\mathcal V_{\mathrm{fin}}=\{S\in\mathcal V: |S|<\infty\},
\qquad
\mathcal V_{\mathrm{fg}}=\{S\in\mathcal V: S=\langle A\rangle\text{ for some finite nonempty }A\subseteq S\}.
\]
Here \(\langle A\rangle\) consists of all finite nonempty products of elements of \(A\).

The supplied question, attributed to G. Mashevitzky, has two indicated versions:
\begin{enumerate}
\item Does there exist a semigroup variety \(\mathcal V\) such that \(\mathcal V_{\mathrm{fin}}\) forms an \emph{elementary inclusive variety}?
\item Does there exist a semigroup variety \(\mathcal V\) such that \(\mathcal V_{\mathrm{fg}}\) forms an \emph{elementary inclusive variety}?
\end{enumerate}
The source uses the wording ``finite (or finitely generated)''; the two versions above retain these alternatives without requiring the same variety to satisfy both. The supplied record does not define \emph{elementary inclusive variety}. Consequently, the target property remains unspecified, and these questions cannot be made fully precise from the supplied evidence alone. In particular, no identification of that phrase with a first-order elementary class or an ordinary equational variety is assumed. No additional nontriviality restriction on \(\mathcal V\) is supplied.
\subsection{Short English statement}
G. Mashevitzky asks whether some semigroup variety has its finite members, or alternatively its finitely generated members, forming an ``elementary inclusive variety.'' The supplied record does not define this target property, so its intended meaning requires clarification.
\subsection{Sources}
\begin{itemize}
\item Local UnsolvedMath problem registry, permanent ID \(800044\), source code \texttt{catalogue entry}, title \emph{The Lyapin's Notebook --- Problem 46}. The supplied registry labels the problem open; this label is not an independently verified mathematical assessment.
\item Bershadsky, Kublanovsky, Mashevitzky, \emph{The Lyapin's Notebook: Unsolved Problems in Semigroup Theory} (2026), Problem 46. Supplied source URL: \url{https://arxiv.org/abs/2604.04763}. The problem statement attributes the question specifically to G. Mashevitzky; the registry's proposer field lists Bershadsky, Kublanovsky, Mashevitzky.
\item Supplied provenance records access on 2026-10-16 and extraction by \texttt{arxiv\_latex}. These are supplied provenance claims, not an assertion that the source was visited or independently verified for this draft.
\end{itemize}
\end{document}
